(1) 令扰动幅值为 \((U,V,N)\)。在波矢 \((k,0)\) 下，线性系统的系数矩阵为
\[
\mathsf A=\begin{pmatrix}
-i\omega&-f&igk\\ f&-i\omega&0\\ ikH&0&-i\omega
\end{pmatrix},\qquad
\det\mathsf A=i\omega(\omega^2-f^2-c^2k^2).
\]
因此三支频率为
\[
\boxed{\omega_0=0,\qquad \omega_\pm=\pm\sqrt{f^2+c^2k^2}}.
\]
对 \(\omega\ne0\)，连续性方程给 \(U=\omega N/(Hk)\)，南北动量方程给
\(V=-ifU/\omega=-ifN/(Hk)\)。非零频率支的横向速度相位相对纵向速度差 \(\pi/2\)，即受科里奥利力影响的椭圆极化。对 \(\omega=0\)，第三式给 \(U=0\)，第一式给 \(V=igkN/f\)；它是静态地转平衡 \(-fv=-g\eta_x\)，而非传播波。直接代回三式可验证其零频率本征矢。惯性重力波在 \(k\to0\) 时仍有非零惯性频率 \(|f|\)。

(2) 设 \(\eta=N(y)e^{i(kx-\omega t)}\)、\(u=U(y)e^{i(kx-\omega t)}\)、\(v=0\)。由沿岸动量及连续性方程，
\[
-i\omega U=-igkN,\qquad -i\omega N+ikHU=0,
\]
从而 \(\omega^2=c^2k^2\)，且 \(U=\omega N/(Hk)\)。横岸动量方程则要求
\[
fU=-gN'(y),\qquad
N(y)=N(0)\exp\!\left[-\frac{f\omega}{c^2k}y\right].
\]
在 \(f>0\)、\(y\ge0\) 的海洋中，衰减条件要求 \(\omega/k>0\)，故只能取
\[
\boxed{\omega_K=ck,\quad u=\frac cH\eta,\quad v=0,\quad
\eta\propto e^{-y/L_R},\quad L_R=\frac cf}.
\]
另一个代数根 \(\omega=-ck\) 会产生 \(e^{+fy/c}\)，不是有限能量的沿岸束缚模。海岸无穿透条件由 \(v=0\) 自动满足；若 \(f<0\) 且海洋仍位于 \(y\ge0\)，允许传播方向反转为 \(\omega=-ck\)，衰减长度为 \(c/|f|\)。该方向选择依赖海岸所在侧，不能仅由无界域的体波频散判断。

(3) 对第二条动量方程取 \(x\) 导数、第一条取 \(y\) 导数并相减，得到
\[
\zeta_t+f(u_x+v_y)=0.
\]
又由连续性方程 \(\eta_t=-H(u_x+v_y)\)，于是
\[
\boxed{\partial_t\left(\zeta-\frac fH\eta\right)=0}.
\]
这是完整位涡 \((f+\zeta)/(H+\eta)\) 在静止背景附近的一阶线性扰动（乘以常数 \(H\) 后）的守恒形式。将两条动量方程分别乘 \(Hu,Hv\)，连续性方程乘 \(g\eta\) 并相加，科里奥利交叉项 \(Hfuv-Hfvu\) 恰好抵消，得到
\[
\boxed{\partial_t\mathcal E+\partial_xF_x+\partial_yF_y=0},\qquad
\mathcal E=\frac H2(u^2+v^2)+\frac g2\eta^2,\quad
(F_x,F_y)=gH\eta(u,v).
\]
因此旋转只改变传播与极化，不直接注入机械能；刚性海岸上 \(F_y=0\)。对 Kelvin 波，\(v=0\)，\(\zeta=-u_y=(f/c)u=(f/H)\eta\)，所以 \(\boxed{q'=0}\)，与其无位涡异常的边界束缚波性质一致。

(4) 正频体波满足 \(\omega_+(k)=\sqrt{f^2+c^2k^2}\)，故
\[
\boxed{v_g=\frac{d\omega_+}{dk}=\frac{c^2k}{\sqrt{f^2+c^2k^2}}}.
\]
当 \(k>0\) 时 \(0<v_g<c\)，而相速 \(\omega_+/k>c\)。长波 \(ck\ll f\) 下
\(\omega_+=f+c^2k^2/(2f)+O(k^4)\)、\(v_g\sim c^2k/f\to0\)；
短波 \(ck\gg f\) 下 \(\omega_+=ck+f^2/(2ck)+O(k^{-3})\)、\(v_g\to c\)。
在固定 \(k>0\) 下 \(f\to0^+\)，体波趋于非旋转浅水重力波 \(\omega=ck\)。
对实 Kelvin 波，由 \(u=(c/H)\eta\) 及 \(c^2=gH\)，逐点有
\[
\mathcal E=g\eta^2,\qquad F_x=gH\eta u=cg\eta^2=c\mathcal E.
\]
记 \(L_R=c/f\)，时间平均 \(\langle\cos^2\rangle=1/2\)，于是
\[
\boxed{\overline{\mathscr E}=\int_0^\infty\langle\mathcal E\rangle\,dy
=\frac{g\eta_0^2}{2}\int_0^\infty e^{-2y/L_R}dy
=\frac{g\eta_0^2L_R}{4}},\qquad
\boxed{\overline{\mathscr F}_x=c\,\overline{\mathscr E}}.
\]
两者比值为 \(c\)，等于 Kelvin 波群速度与相速。对固定岸边振幅 \(\eta_0\)，\(f\to0^+\) 时 \(L_R\to\infty\)，积分能量发散，表明横向束缚态不再可归一化；不能把这一极限解释成有限宽度的沿岸波。单位密度下 \([\overline{\mathscr E}]=\mathrm{m}^4\mathrm{s}^{-2}\)、\([\overline{\mathscr F}_x]=\mathrm{m}^5\mathrm{s}^{-3}\)，二者之比确为速度。
